\section{Introduction}\label{sec:introduction}

% Paragraph plan: notes/outline.md, Section 1. The comments mark the role of each paragraph.

% (1) Background: industrial practice and the trade-off between cost and model use.
Feedback control has developed rapidly over the past decades, yet the algorithms deployed at
scale in industrial production remain few.
Proportional--integral--derivative (PID) control, the linear quadratic regulator (LQR), and
model predictive control (MPC) are the most representative of them: in a survey conducted by
the IFAC Industry Committee, PID control received the highest rating for current industrial
impact and MPC the highest rating for expected future impact~\citep{samad2020industry}.
These algorithms occupy different positions in a trade-off between computational cost and
the ability to exploit a model of the plant, that is, to anticipate its motion and to respect
its constraints.
Many emerging applications, such as agile robots, automated vehicles, and aerial vehicles,
demand both at once: their dynamics are nonlinear and often underactuated, and their
controllers must run at high rates on embedded processors with limited computing resources.

% (2) PID: model-free and cheap, but reactive and sensitive to its error history.
PID control is the representative of model-free feedback.
Its control signal combines the present error, the accumulated past error, and the rate of
change of the error, its tuning requires no model of the plant, and it remains the workhorse
of industrial feedback~\citep{astrom2001future}.
Its computation amounts to a few arithmetic operations per sample, so that it runs comfortably
on microcontrollers.
These advantages come with limitations.
Because the integral term stores the entire error history, past errors keep acting on the
plant long after they have occurred: transient disturbances and measurement noise that enter
the integrator are never forgotten, and while the actuator saturates, the integrator keeps
accumulating error that is released only after the actuator desaturates.
The result can be large overshoot and sustained oscillation, a phenomenon known as integrator
windup~\citep{astrom1989windup,astrom2008feedback}.
The derivative term, in turn, amplifies high-frequency measurement noise and must be
filtered~\citep{astrom2008feedback}.
Moreover, PID control is essentially reactive.
Its only anticipation is the linear extrapolation of the error implied by the derivative
term~\citep{astrom2008feedback}, so it can anticipate neither the nonlinear motion of the plant
nor upcoming changes of the reference, and its tracking performance deteriorates when the
reference varies rapidly.
Finally, although PID control appears conceptually simple, tuning it to meet several
conflicting objectives is hard in practice~\citep{ang2005pid}, and the difficulty is
compounded in strongly coupled multivariable and underactuated systems, where the loops are
designed channel by channel.

% (3) LQR: model-based and optimal, but Riccati equations, no constraints, local only.
The LQR is the prototype of model-based optimal
feedback~\citep{kalman1960contributions,anderson1990optimal}.
For a linear model and a quadratic cost, the optimal control is a linear state feedback whose
gain follows from a Riccati equation.
For linear time-invariant regulation, this equation is solved once, offline.
For nonlinear plants that operate over wide ranges or track time-varying references, however,
the model must be linearized along the operating trajectory and the Riccati equation solved
again online, or the gains must be scheduled~\citep{rugh2000research}, which raises the
computational burden markedly above that of PID control.
Standard LQR cannot enforce input or state constraints: saturating its output forfeits its
optimality and stability guarantees, and incorporating the constraints turns the control law
into a piecewise affine function of the state that must be computed by multiparametric
programming~\citep{bemporad2002explicit}.
LQR has further limitations.
When it is applied to a nonlinear plant through linearization, stability is guaranteed only in
a neighborhood of the operating point, so that large-amplitude maneuvers such as the swing-up
of a pendulum lie outside its scope.
The weighting matrices of the cost are only indirectly related to time-domain specifications
and are usually tuned by trial and error.
And the celebrated robustness margins of LQR~\citep{safonov1977gain} hold only under full-state
feedback: no margins are guaranteed once a state observer is inserted into the
loop~\citep{doyle1978guaranteed}.

% (4) MPC: constraints and nonlinear models, but an optimization problem per sample.
MPC overcomes several of these limitations by solving, at every sampling instant, a
finite-horizon optimal control problem for the current state and applying the first element of
the optimal input sequence~\citep{mayne2000constrained}.
It handles multivariable and nonlinear models, incorporates input and state constraints
explicitly, and has become the standard advanced control technology in the process
industries~\citep{qin2003survey}.
Its price is computational: a quadratic program, or a nonlinear program for nonlinear models,
must be solved within each sampling period.
Efficient solvers exploit the sparse, stage-wise structure of these problems through structured
or sparse matrix factorizations~\citep{rao1998application,wang2010fast,stellato2020osqp}, and
real-time schemes either perform only a single Newton-type iteration per
sample~\citep{diehl2005realtime,verschueren2022acados} or, as continuation methods do, update
the input sequence by integrating a differential equation that drives the residual of the
optimality conditions to zero~\citep{ohtsuka2004continuation}; nevertheless, the computational
and memory demands remain far above those of LQR.
Explicit MPC moves the optimization offline, but the number of regions of the resulting
piecewise affine law grows rapidly with the size of the problem~\citep{bemporad2002explicit}.
MPC therefore remains difficult to deploy on low-cost embedded platforms with fast sampling.
In addition, its closed-loop stability rests on suitably designed terminal costs and
constraints~\citep{mayne2000constrained}.

% (5) iLQR: nonlinear extension of LQR, at a higher computational cost.
The iterative linear quadratic regulator (iLQR)~\citep{li2004iterative} extends LQR to
nonlinear dynamics.
It is a variant of differential dynamic programming~\citep{mayne1966second} that neglects the
second-order derivatives of the dynamics: each iteration linearizes the dynamics around the
current trajectory, solves the resulting time-varying LQR problem by a backward Riccati-type
recursion, and updates the trajectory by a forward rollout.
Applied in a receding-horizon fashion, iLQR-type methods have synthesized complex behaviors
online, although for a humanoid robot the computation still ran about seven times slower than
real time on a standard PC~\citep{tassa2012synthesis}; constrained variants have been adopted
for motion planning in autonomous driving~\citep{chen2017constrained,chen2019autonomous}.
Every control update requires repeated forward--backward passes over the horizon, so that the
computational cost exceeds that of LQR; constraints are not handled natively and call for
dedicated extensions~\citep{tassa2014control,chen2019autonomous}; and, being a local method,
iLQR may converge to poor local solutions depending on its initialization.

% (6) Structured nonlinear designs: rigorous but system-specific.
Beyond these industrial workhorses, nonlinear control theory offers constructive designs with
rigorous guarantees, including feedback linearization~\citep{isidori1995nonlinear},
backstepping~\citep{krstic1995nonlinear}, sliding-mode control~\citep{utkin1977variable}, and
passivity-based control~\citep{ortega2002interconnection}; see \citet{khalil2002nonlinear} for
an overview.
These designs exploit structural properties of the model, for instance strict-feedback forms in
backstepping or the solvability of matching partial differential equations in the energy
shaping of underactuated mechanical systems~\citep{ortega2002stabilization}, and they are
largely system-specific.
Correspondingly, industrial surveys rate the current impact of nonlinear control far below that
of PID control and MPC~\citep{samad2020industry}.

% (7) Research gap.
The picture that emerges is one of two extremes.
PID control requires negligible computation but neither uses a model nor anticipates the motion
of the plant, whereas LQR, iLQR, and MPC exploit a model at a computational cost that grows from
solving Riccati equations to solving an optimization problem at every sampling instant.
What is missing is a controller that exploits a nonlinear model to anticipate the motion of the
plant, acts on every state component, including the unactuated states of underactuated systems,
keeps the per-step computation low enough for low-cost embedded hardware, and comes with
convergence guarantees.
\TODO{研究空白；“单步计算量低”需在实验中以实测耗时支撑（笔记 R11）}

% (8) Speed-gradient method: the starting point of HFC.
The speed-gradient (SG) method introduced by Fradkov~\citep{fradkov1979speed,fradkov1999nonlinear}
points in this direction.
Consider a plant $\dot\bx = \mathbf{f}(\bx,\bu,t)$ and a goal function $Q(\bx,t) \ge 0$ whose
vanishing encodes the control objective.
The SG method computes the rate of change, or \emph{speed}, of $Q$ along the trajectories of the
plant, $\omega(\bx,\bu,t) = \partial Q(\bx,t)/\partial t + \nabla_{\bx} Q(\bx,t)^\T
\mathbf{f}(\bx,\bu,t)$, and changes the input against the gradient of this speed with respect to
the input: in the differential form $\dot\bu = -K \nabla_{\bu}\omega(\bx,\bu,t)$ or, for plants
affine in the input, in the finite form $\bu = -K \nabla_{\bu}\omega(\bx,\bu,t)$, where $K$ is a
positive definite gain matrix.
Under a convexity condition on $\omega$ with respect to $\bu$ and an attainability condition, which
requires an input that renders $\omega$ sufficiently negative, the goal $Q \to 0$ is achieved,
with the goal function serving as a Lyapunov function~\citep{fradkov1999nonlinear}.
The method is constructive and model-based; its computation per step reduces to the evaluation
of a gradient, with neither a Riccati equation nor an online optimization; and it applies to
nonlinear, non-affine, and time-varying systems.
It has accordingly served as a unifying framework for adaptive control~\citep{fradkov1979speed},
energy-based control of oscillations~\citep{fradkov1996swinging}, stabilization of invariant
sets~\citep{shiriaev2000stabilization}, and control of chaos~\citep{fradkov2007cybernetical}.
Its limitations stem from its reliance on the instantaneous speed of the goal function.
First, $\nabla_{\bu}\omega$ receives no contribution from the state components whose rates do not
depend on the input.
In mechanical and kinematic models, the input acts directly only on the highest-order states,
such as velocities, so that errors in the low-order states, such as positions, drop out of the
control law.
The designer must then construct an auxiliary goal function, typically energy-like, whose speed
depends explicitly on the input; since the attainability condition essentially requires this
function to be a control Lyapunov function, the construction is difficult for general
underactuated systems.
Second, the resulting control is myopic: it decreases $Q$ at the current instant without
anticipating the future motion of the plant.
Third, achieving $Q \to 0$ does not in general imply convergence of the full state.
Energy-based swing-up laws, including those obtained by the SG
method~\citep{fradkov1996swinging}, regulate the energy of a pendulum to that of its upright
equilibrium and thereby steer the state toward the corresponding homoclinic orbit rather than to
the equilibrium itself, so that a separate local stabilizer must be switched in near the upright
position~\citep{astrom2000swinging,shiriaev2000global}.

% (9) Flow-based cost of Lan and Cvitanovic: the second ingredient of HFC.
A different line of work, originating in the computation of unstable periodic orbits, shows how a
cost can be built from the flow itself.
\citet{lan2004variational} proposed a variational method in which a trial loop in phase space is
deformed toward a periodic orbit.
Their cost measures, along the loop, the mismatch $\lambda\bv - \tilde{\bv}$ between the velocity
field $\bv$ of the flow, rescaled by a local time-scaling factor $\lambda$, and the tangent
$\tilde{\bv}$ of the loop; in discretized form, it compares the image of each loop point under
the flow over one time step with the next loop point.
The loop evolves under a ``Newton descent'' equation, an infinitesimal-step version of the damped
Newton--Raphson iteration, along which this mismatch decays exponentially in a fictitious time,
so that the cost decreases monotonically.
For systems with higher time derivatives, the method matches the highest derivatives instead.
The idea of measuring the cost through the phase-space flow, rather than through a hand-crafted
function of the state, is the one we adopt in this paper.

% (10) Controllers whose input evolves as a flow: closest related work.
Letting the control input evolve as the solution of a differential equation that drives a
residual to zero is not new in control either.
Besides the continuation methods for receding-horizon control mentioned
above~\citep{ohtsuka2004continuation}, Wardi and co-workers have developed an output-tracking
controller based on a fluid-flow version of the Newton--Raphson method: an integrator with variable
gain that drives to zero the difference between a predicted future output and the reference
signal~\citep{wardi2017performance,wardi2024tracking}.
Implementable by simple algorithms, it has been applied to autonomous
vehicles~\citep{shivam2019tracking}, to aggressive quadrotor flight, where it outperformed the native
tracking controller of the PX4 autopilot~\citep{morales2024newton}, and to computationally
constrained aerial vehicles, where it matched or exceeded nonlinear MPC at a substantially lower
computation time~\citep{morales2026lightweight}; stability conditions are available for linear
systems and for a class of differentially flat systems~\citep{wardi2024tracking,niu2025stability}.
\method{} shares with this controller the idea of steering a predicted state toward the reference by
a flow of the input, but it differs in two respects.
First, the Newton--Raphson flow solves a square system of equations, with as many tracked outputs as
inputs, and its update requires the inverse of the Jacobian of the predicted output with respect to
the input.
\method{} instead weighs the predicted error of every state component and drives the input along a
gradient scaled by diagonal coefficients, so that no matrix is inverted and the states outside the
tracked outputs, such as the pendulum angles of an underactuated system, enter the control law
directly.
Second, in its basic form the Newton--Raphson flow predicts the output by integrating the plant over
the preview horizon with the input held constant~\citep{wardi2024tracking}, whereas \method{}
evaluates a Taylor polynomial of the flow in closed form, whose higher-order terms already make every
state component depend on the input.

% (11) The proposed approach.
In this paper, we propose \methodlong{} (\method{}), which builds on the SG method and on the
flow-based cost of \citet{lan2004variational}.
The reference is a sequence of state samples without time stamps, whose evolution times are close
to the update interval $h$ of the controller but not exactly uniform; at every round, the controller
takes the $m$th next sample as target, which requires no projection onto the reference, and aims to
reach it after the preview time $\tau = mh$.
Writing $\bv_\tau$ for the mean velocity of a look-ahead of the plant over the preview time, the cost
\[
    F(\bx,\bu,t) = \norm[\big]{\Gamma\bigl(\tau\bv_\tau(\bx,\bu) - (\bxref - \bx)\bigr)}^2
\]
compares the flow of the plant with the mean velocity that would carry the current state to the
target sample, the matrix $\Gamma$ rendering the cost dimensionless.
This is the counterpart of the mismatch $\lambda\bv - \tilde{\bv}$ of \citet{lan2004variational}, the
secant toward the target sample playing the role of the loop tangent and the preview time that of
the time-scaling factor: the controller chases a flow rather than a geometric position.
The control input evolves along the negative gradient of $F$ with respect to $\bu$, scaled by
dimension-balancing coefficients, and the controller follows this gradient flow by one explicit
update per round, without any online optimization.
To first order in $\tau$, this control law coincides, up to the gain matrix, with the SG law for the
quadratic goal function $Q = \frac{1}{2}\norm{\Gamma(\bx - \bxref)}^2$ (\cref{prop:sg}), and a
look-ahead of first order would therefore inherit the first limitation of SG: the low-order states
would not contribute to the gradient.
\method{} removes this limitation by taking as look-ahead the Taylor polynomial of the flow of the
plant, with the input held over the preview time, up to the order at which the input reaches every
state component: the low-order states are predicted through the time derivatives of the high-order
ones, much as \citet{lan2004variational} match the highest derivatives of systems of higher order,
so that every state component contributes to the control update.
As a result, and unlike SG, \method{} needs no hand-crafted energy-like goal function: a quadratic
cost on the state errors suffices.
Expanding the square, $F/2 = Q + \tau\mathcal{H}$, where $\mathcal{H}$ has the form of the Hamiltonian
of Pontryagin's principle~\citep{pontryagin1962mathematical}, with the gradient of $Q$ in the role of
the costate, the mean velocity of the look-ahead in that of the velocity, and the kinetic energy of
this velocity as running cost (\cref{rem:hamiltonian}).
The input therefore descends a control Hamiltonian continuously, rather than minimizing it along the
solution of an adjoint equation, which motivates the name of the method.
We call $F$ the flow cost; as the sum of the potential $Q$ and $\tau$ times this Hamiltonian, it is an
energy-like function, and \method{} is an energy-descent control law.
The Hamiltonian is not meant in the sense of Hamiltonian mechanics: the closed loop is dissipative,
and no symplectic structure is claimed.
The controller uses the target samples only: neither their timing, nor values between them, nor
the input that generates them is required.

% (12) Contributions.
The main contributions of this paper are as follows.
\begin{itemize}
    \item We propose \method{}, which drives the control input along the gradient of a cost that
    compares the flows of the plant and of the reference over a preview time, in the spirit of the
    variational method of \citet{lan2004variational}; we show that this cost is the weighted state
    error plus the preview time times a control Hamiltonian of Pontryagin type, and that the control
    law contains the SG law as its first-order approximation in the preview time.
    \item We introduce a Taylor look-ahead, through which the low-order states, which the input
    reaches only through the higher-order ones, contribute to the control update, so that no
    hand-crafted energy-like goal function is needed; the look-ahead is affine in the input, so
    that the cost is quadratic in it.
    \item We show that on integrator chains the closed-loop eigenvalues of \method{} are $1/\tau$
    times the roots of a polynomial fixed by the weights, which yields exact conditions for
    stability, met by equal weights up to chains of four levels, and for the update interval; we
    show that the look-ahead toward a target sample causes a time delay equal to the mismatch
    between the preview time and the time to the sample, with no offset up to terms of third order
    in the preview time or of second order in the rate of the reference input, that taking the $m$th
    next sample with $\tau = mh$ removes the delay without any projection onto the reference, and we
    bound the resulting tracking error.
    \TODO{一般非线性系统的收敛性定理（类型待定，见 notes/convergence.md）}
    \item We derive \method{} laws for three benchmark problems: the swing-up and stabilization of a
    double inverted pendulum on a cart~\citep{graichen2007swingup} by a single control law
    without switching, trajectory tracking for a kinematic bicycle model, for which we prove local
    exponential stability of straight references at every speed in continuous time and quantify the
    speed range of the sampled-data implementation, and trajectory tracking for a
    quadrotor whose attitude evolves on $\mathrm{SO}(3)$~\citep{lee2010geometric}.
    \item We validate \method{} in simulations and physical experiments on all three platforms.
    \TODO{对比基线与评价指标待定；建议包括 SG 与 Newton–Raphson 流方法（笔记 R6）}
\end{itemize}

% (13) Organization. Draft -- revisit once Sections 2--4 are settled.
The remainder of this paper is organized as follows.
\Cref{sec:formulation} recalls the required geometric preliminaries, introduces the flow cost and
its Taylor look-ahead, defines the control law and its sampled-data implementation, relates the
cost to a control Hamiltonian and the law to the SG method, analyzes the convergence and the
tracking error caused by the look-ahead, and treats non-uniform sample intervals; it then derives
the control laws for the three benchmark problems.
\Cref{sec:experiments} reports simulation and experimental results for these systems, and
\cref{sec:conclusion} concludes the paper and outlines possible extensions.

% (14) Notation.
\paragraph{Notation.}
\TODO{符号说明：待 Formulation 定稿后补充}
